% preamble.tex

\usepackage[utf8]{inputenc}
\usepackage[T1]{fontenc}
\usepackage[english]{babel} % Ho cambiato in italian, supponendo che la maggior parte del testo sia in italiano
\usepackage{amsmath}
\usepackage{amssymb}
\usepackage{bm} % For bold math symbols (\bm command)
\usepackage{graphicx}
\usepackage{hyperref}
\usepackage{geometry}
% \usepackage{fontspec} % Utile per Xetex/LuaLaTeX ma mantenuto - COMMENTED OUT: only works with XeLaTeX/LuaLaTeX
\usepackage{booktabs}
\usepackage{amsthm}
\usepackage{enumitem}
\usepackage{caption}
\usepackage{thmtools}
\usepackage[skins]{tcolorbox} % For colored boxes
\usepackage{standalone} % Già presente nel tuo elenco
\usepackage{tkz-euclide} % Questo carica TikZ
\usetikzlibrary{fit} % Necessaria per l'opzione fit
\usetikzlibrary{shapes.geometric}
\usetikzlibrary{decorations.pathmorphing} % For snake decoration used in TikZ
\usepackage{graphicx}
\usepackage[T1]{fontenc}
% --- PACCHETTI AGGIUNTI PER LA FORMATTAZIONE ---
\usepackage{parskip}  % Aggiunge spazio tra i paragrafi (stile "blocco")
\usepackage{fancyhdr} % Per intestazioni e piè di pagina personalizzati
\usepackage{fontawesome5}
\usepackage{xcolor}       % Per colorare le icone
\usepackage{setspace}
\usepackage[utf8]{inputenc}
\usepackage{amsmath}
\usepackage{amsfonts} % For blackboard bold and other math fonts
\usepackage{minted}   % For colored C++ code (requires Pygments and -shell-escape)
\setminted{
    fontsize=\small, % Makes the font slightly smaller
}
\usepackage{etoolbox} % For hooks to modify environments
\BeforeBeginEnvironment{minted}{\begin{center}}
\AfterEndEnvironment{minted}{\end{center}}
\usepackage{amssymb}
\usepackage{graphicx}
\usepackage{booktabs} % For professional tables
\usepackage{tabularx}
\renewcommand{\arraystretch}{1.4} % Aumenta la spaziatura verticale tra le righe delle tabelle


% --- IMPOSTAZIONI GEOMETRIA E FONT ---
\geometry{a4paper, margin=1in}
%\onehalfspacing
% Configura lo stile 'plain' (usato da \maketitle e \tableofcontents)
% per essere coerente con il piè di pagina
\fancypagestyle{plain}{
    \fancyhf{} % Pulisce
    \cfoot{\thepage} % Mette solo il numero di pagina
    \renewcommand{\headrulewidth}{0pt} % Niente linea su queste pagine
}

\usepackage{fontspec}
\usepackage{unicode-math} % Necessario per xelatex/lualatex
\setmainfont{Fira Sans}[
    % Imposta la variante "Light" come la variante principale (Regular/Medium)
    % Quando usi \setmainfont senza specificare 'Path', cerca il nome completo del font.
    UprightFont = Fira Sans Light,
    % Quando usi \textbf{}, userà questa variante (Bold)
    BoldFont = Fira Sans Regular, 
    % Quando usi \textit{}, userà questa variante (Light Italic)
    %ItalicFont = Fira Sans Light Italic,
    % Quando usi \textbf\textit{}, userà questa variante (Bold Italic)
    %BoldItalicFont = Fira Sans Bold Italic,
]
\setsansfont{Fira Sans}
\setmathfont{Libertinus Math}

% Configurazione Header/Footer per le altre pagine
\pagestyle{fancy}
\fancyhead{} % Pulisce l'intestazione
\fancyfoot{} % Pulisce il piè di pagina
\fancyhead[L]{\nouppercase{\rightmark}} % Nome Sezione a sinistra (o Chapter)
\fancyhead[R]{\nouppercase{\leftmark}}  % Nome Capitolo a destra (o Sezione)
\fancyfoot[C]{\thepage} % Numero di pagina al centro
\renewcommand{\headrulewidth}{0.4pt}
\renewcommand{\footrulewidth}{0pt}



% Definizioni utili per il Main file
\newcommand{\inputlesson}[1]{\clearpage\input{#1}\clearpage} % Comando per input lezione

